\section{The Congestion-Stress Paradox in Bottleneck Data}

Paper A.1 identified the congestion-stress paradox in the aggregate score distribution: urban beltways score higher than transcontinental arteries because aggregate scores weight throughput gap heavily, and short high-AADT corridors accumulate maximum A1 scores. The ATRI bottleneck data provides an independent validation of this paradox.

\subsection{Tier Misclassification in ATRI Data}

If the aggregate-score tier classification were correct, T1 corridors would account for a proportionate share of ATRI bottleneck locations. The aggregate-score T1 tier (13 corridors, including I-110, I-880, I-285, I-4) would be expected to dominate the ATRI rankings because these corridors score highest on congestion metrics.

In fact, 3 of the 13 aggregate-score T1 corridors (I-285, I-4, I-880) account for only 10 ATRI top-50 locations and \$4.6B in annual cost --- a lower per-location cost than the centrality-adjusted T1 corridors (I-75, I-95, I-80, I-10).

The centrality-adjusted T1 corridors (8 corridors: I-5, I-10, I-35, I-40, I-75, I-80, I-90, I-95) account for 18 ATRI top-50 locations and \$14.1B in annual cost. The cascade multiplier confirms that centrality-adjusted T1 bottlenecks are more economically significant than aggregate-score T1 bottlenecks.

\subsection{The Atlanta Paradox}

Atlanta presents the clearest illustration. The I-285 beltway (aggregate-score T1, centrality-adjusted T2) has 4 of the top 50 ATRI locations (\$3.0B/yr). I-75 Atlanta segment (aggregate-score T2, centrality-adjusted T1) has 2 ATRI locations but \$1.3B/yr. Per location, I-75 costs more (\$650M vs. \$750M for I-285) --- not dramatically different, but the I-75 locations cascade northward through the Great Lakes manufacturing corridor in a way that the I-285 locations do not.

The investment conclusion is the same as in A.1: T2 relief (improving I-285 throughput to reduce the number of trucks funneling through the interchange) provides more bottleneck locations per dollar of investment, but T1 managed lanes (improving I-75 throughput) provides more total economic cost reduction per dollar.

\subsection{The I-40 Exception Confirmed}

I-40 has zero ATRI top-50 locations. This is the paper's strongest validation of the tier classification: the one T1 corridor where our analysis predicted no capacity emergency (V/C = 0.84, at target) is also the one T1 corridor with no ATRI-ranked bottlenecks. The convergence between ROUTE's scoring model and ATRI's revealed-preference data for I-40 significantly strengthens the credibility of the classification for the other 7 T1 corridors.
